% © 2026 Ing. David Jaroš — CC BY-NC-ND 4.0
%
% This work is licensed under a Creative Commons Attribution-NonCommercial-NoDerivatives
% 4.0 International License (CC BY-NC-ND 4.0).
%
% License History: Earlier drafts (up to v0.3) were released under CC BY 4.0.
% From v0.4 onward, all material is released under CC BY-NC-ND 4.0 to protect
% the integrity of the theoretical work during ongoing academic development.
%
% See LICENSE.md for full license text.
%
% File: research_tracks/alpha_spectral/b_coefficient_gap_resolution.tex
%
% Task: close_or_kill_B_coefficient_gap
% Priority: CRITICAL
% Gap ID: G137-B
%
% Purpose: Determine whether B(p) = (p+1)/3 can be derived from the UBT
%   action S[Theta] or must remain a phenomenological/modular ansatz.
%   Synthesises all existing analysis (RG, KK+winding, modular index,
%   Chowla-Selberg, fixed-point bootstrap) and issues a definitive verdict.
%
% Hard rules:
%   - No use of observed alpha as input
%   - No fitting B to 137
%   - If B cannot be derived, state CONDITIONAL clearly

\documentclass[12pt,a4paper]{article}
\usepackage[a4paper, margin=2.5cm]{geometry}
\usepackage{amsmath, amssymb, amsthm}
\usepackage{mathtools}
\usepackage{hyperref}
\usepackage{xcolor}
\usepackage{booktabs}
\usepackage{longtable}
\usepackage{mdframed}
\usepackage{tcolorbox}
\tcbuselibrary{skins,breakable}

\newtheorem{theorem}{Theorem}[section]
\newtheorem{lemma}[theorem]{Lemma}
\newtheorem{proposition}[theorem]{Proposition}
\newtheorem{corollary}[theorem]{Corollary}
\theoremstyle{definition}
\newtheorem{definition}[theorem]{Definition}
\newtheorem{gap}{Open Gap}[section]
\theoremstyle{remark}
\newtheorem{remark}[theorem]{Remark}

\newmdenv[backgroundcolor=yellow!15, linecolor=orange!80, linewidth=1.5pt]{gapbox}
\newmdenv[backgroundcolor=green!10, linecolor=green!60!black, linewidth=1.5pt]{resultbox}
\newmdenv[backgroundcolor=red!8,   linecolor=red!55,   linewidth=1.5pt]{warnbox}
\newmdenv[backgroundcolor=blue!6,  linecolor=blue!50,  linewidth=1.5pt]{insightbox}
\newmdenv[backgroundcolor=gray!8,  linecolor=gray!50,  linewidth=1.0pt]{infobox}

\newcommand{\proved}{\textcolor{green!55!black}{\textbf{[PROVED]}}}
\newcommand{\heur}{\textcolor{orange!80!black}{\textbf{[HEURISTIC]}}}
\newcommand{\cond}{\textcolor{blue!70!black}{\textbf{[COND]}}}
\newcommand{\openp}{\textcolor{red!70!black}{\textbf{[OPEN]}}}
\newcommand{\obs}{\textcolor{blue!50!black}{\textbf{[OBS]}}}

\title{\textbf{B-Coefficient Gap Resolution}\\[4pt]
       \large Can $B(p)=(p+1)/3$ Be Derived from the UBT Action $S[\Theta]$?\\[2pt]
       \normalsize Gap G137-B --- Final Analysis}
\author{Ing.\ David Jaro\v{s}\\
  \small \texttt{research\_tracks/alpha\_spectral/}}
\date{May 2026}

\begin{document}
\maketitle

\begin{abstract}
We investigate whether the coefficient $B(p) = (p+1)/3$ appearing in the UBT
effective potential $V_{\mathrm{eff}}(n) = n^2 - B\,n\ln n$ --- a coefficient
required to place the prime attractor at $n^* = p$ --- can be derived
from the UBT action $S[\Theta]$ without using experimental inputs.
Four routes are examined: (1) one-loop renormalisation group (RG) from the
Kaluza--Klein and winding spectrum; (2) the KK+winding estimate yielding
$B \approx 43.6$; (3) the modular-index route $B = \mu(\Gamma_0(p))/3$; and
(4) a fixed-point bootstrap self-consistency argument.
We assess gauge invariance and regulator dependence of each candidate.
\textbf{Verdict}: Routes (1) and (2) are incomplete by a factor
$\approx 1.054$; route (3) produces the correct numerical value but cannot
currently be derived from $S[\Theta]$; route (4) confirms self-consistency
but does not constitute a derivation.
\textbf{The result $B(p)=(p+1)/3$ must remain a \emph{CONDITIONAL} modular ansatz
until the missing mechanism is identified.}
No experimental value of $\alpha$ was used in any step.
\end{abstract}

\tableofcontents
\newpage

%────────────────────────────────────────────────────────────────────────────
\section{Setup and Notation}
\label{sec:setup}
%────────────────────────────────────────────────────────────────────────────

\subsection{The effective potential and the role of $B$}

The UBT winding-mode effective potential on the compact imaginary-time circle
$S^1_\psi$ of radius $R_\psi$ is
\begin{equation}
  V_{\mathrm{eff}}(n) = n^2 - B\,n\ln n,
  \qquad n \in \mathbb{Z}_{>0}.
  \label{eq:veff}
\end{equation}
The minimum of~\eqref{eq:veff} occurs at $n^* = n^*(B)$ satisfying the
fixed-point equation
\begin{equation}
  2n^* = B\bigl(\ln n^* + 1\bigr).
  \label{eq:fixedpoint}
\end{equation}
For the prime attractor to coincide with the prime $p = 137$, we need
\begin{equation}
  B_{\mathrm{phenom}}
  = \frac{2 \times 137}{\ln 137 + 1}
  = \frac{274}{4.9189 + 1}
  \approx 46.298.
  \label{eq:Bphenom}
\end{equation}

\subsection{The phenomenological formula}

The prime-stability analysis (see \texttt{canonical/alpha/prime\_stability\_set.tex})
shows that stable primes (prime attractors of $V_{\mathrm{eff}}$) arise when
\begin{equation}
  B(p) = \frac{p+1}{3}.
  \label{eq:Bformula}
\end{equation}
At $p = 137$: $B(137) = 138/3 = 46$.  The value~\eqref{eq:Bformula} at
$p = 137$ agrees with $B_{\mathrm{phenom}}$~\eqref{eq:Bphenom} to within 0.64\%
(discrepancy $\Delta B \approx 0.298$).

\begin{warnbox}
\textbf{Hard rule}: Equation~\eqref{eq:Bformula} is an ansatz to be examined,
not an axiom.  No experimental value of $\alpha$ or any observed physical
constant enters any step.
\end{warnbox}

\subsection{Known proved steps}

The following steps are established at level [L0] or [L1] and are not
re-examined here:

\begin{enumerate}
  \item $N_{\mathrm{eff}} = 12$ from the real dimension of $\mathbb{C}\otimes\mathbb{H}$ [L0].
  \item KK mass spectrum $m_n^2 = n^2$ in natural units [L0].
  \item One-loop $\delta V \propto n^2\ln n$ in $d=1$ [L1].
  \item Gauge invariance of $V_{\mathrm{eff}}$ ($n$ is a $\mathrm{U}(1)_\psi$ charge) [proved].
  \item Regulator independence of the $n\ln n$ coefficient [proved].
  \item $n^*(B_{\mathrm{phenom}}) = 137$ (given $B = B_{\mathrm{phenom}}$) [proved numerically].
\end{enumerate}

\noindent The single remaining gap is the derivation of $B = B_{\mathrm{phenom}}$
from $S[\Theta]$ (Gap G137-B, see \texttt{reports/alpha\_missing\_lemma.md}).

%────────────────────────────────────────────────────────────────────────────
\section{Route 1: One-Loop RG from KK Spectrum}
\label{sec:rg}
%────────────────────────────────────────────────────────────────────────────

\subsection{One-loop calculation}

The one-loop Coleman--Weinberg effective potential in $d=1$ for a scalar with
mass $m_n^2 = n^2$ is (dimensional regularisation, $\overline{\mathrm{MS}}$):
\begin{equation}
  \delta V_{\mathrm{1-loop}}(n)
  = -\frac{n^2}{4\pi}\!\left(\ln\frac{n^2}{\mu^2} - 1\right)
  = -\frac{n^2}{2\pi}\ln n + \text{($\mu$-dep.\ and const.\ terms)}.
  \label{eq:dV1loop}
\end{equation}
After absorbing the $\mu$-dependent term into the renormalisation of the $n^2$
coefficient:
\begin{equation}
  V_{\mathrm{ren}}(n) = n^2 - \frac{n}{2\pi}\cdot n\ln n,
  \label{eq:Vren1loop}
\end{equation}
which corresponds to $B_{\mathrm{1-loop}}(n) = n/(2\pi)$ --- an $n$-dependent
quantity.  At $n = 137$: $B_{\mathrm{1-loop}}(137) \approx 137/(2\pi) \approx 21.8$.

\begin{gapbox}
\textbf{Gap 1a (n-dependent B)}: The naive one-loop result yields an
$n$-dependent coefficient $B(n) = n/(2\pi)$, not a constant.  A constant
$B$ requires either (i) evaluating at a fixed $n = n^*$, introducing a
self-reference, or (ii) an additional physical mechanism.
\end{gapbox}

\subsection{RG consistency check}

The Callan--Symanzik equation for $V_{\mathrm{eff}}(n;\mu)$ is satisfied at
one loop by~\eqref{eq:dV1loop} with the standard $d=1$ beta function.
The $\ln n$ coefficient is indeed scheme-independent (proved in
\texttt{research\_tracks/rg\_nlogn/full\_rg\_derivation.tex}), confirming
that the $n\ln n$ structure is physically meaningful.

\paragraph{Regulator independence.}
\proved\ The coefficient of $n^2\ln(n^2/\mu^2)$ in dimensional regularisation
is UV-finite and scheme-independent at one loop.  Regulator dependence appears
only in the $n^2\ln\mu^2$ term, which is absorbed into the $n^2$ renormalisation.

\paragraph{Gauge invariance.}
\proved\ The effective potential~\eqref{eq:veff} depends only on the integer KK
level $n$, which is the charge under the compact translation symmetry
$\mathrm{U}(1)_\psi$ on $S^1_{L_\psi}$.  This is gauge-invariant under the
$\mathrm{SU}(3)\times\mathrm{SU}(2)\times\mathrm{U}(1)$ interaction sector.

\subsection{Conclusion of Route 1}

\begin{resultbox}
Route 1 gives: $B_{\mathrm{1-loop}} \approx 21.8$ for $n = 137$.
Status: \heur\ (correct structure, wrong magnitude by factor $\approx 2.1$).
The $n^2\ln n$ shape is rigorous; the coefficient $B = 46$ is not reproduced.
\end{resultbox}

%────────────────────────────────────────────────────────────────────────────
\section{Route 2: KK+Winding Estimate at Self-Dual Radius}
\label{sec:kk_winding}
%────────────────────────────────────────────────────────────────────────────

\subsection{T-duality degeneracy}

At the self-dual radius $R_\psi = 1$ (T-duality fixed point), the winding-mode
masses $m_{\mathrm{wind}}^2 = n^2 R_\psi^2 = n^2$ coincide with the KK masses
$m_{\mathrm{KK}}^2 = n^2/R_\psi^2 = n^2$.  The combined one-loop vacuum energy
receives a factor of 2 from the KK/winding degeneracy:
\begin{equation}
  B_{\mathrm{KK+wind}} = 2 \times B_{\mathrm{1-loop}} = 2 \times 21.8 = 43.6.
  \label{eq:Bkkwind}
\end{equation}

\heur\ The factor of 2 from T-duality degeneracy requires:
\begin{enumerate}
  \item The self-dual radius $R_\psi = 1$ is not derived from UBT moduli; it
        is assumed.
  \item The two contributions (KK and winding) must add coherently (not cancel);
        this follows from the sign of the one-loop vacuum energy in $d=1$ but has
        not been verified for the full UBT field content.
\end{enumerate}

\subsection{The missing gap: $43.6 \to 46$}

\begin{gapbox}
\textbf{Gap 2 ($43.6 \to 46$ discrepancy)}: The KK+winding estimate gives
$B \approx 43.6$, which is $\Delta B \approx 2.4$ below the required value
$B_{\mathrm{phenom}} \approx 46.298$.  The relative discrepancy is
\[
  \frac{B_{\mathrm{phenom}} - B_{\mathrm{KK+wind}}}{B_{\mathrm{phenom}}}
  \approx \frac{2.7}{46.298} \approx 5.8\%.
\]

\textbf{Candidate sources of $\Delta B \approx 2.4$}:
\begin{enumerate}
  \item \textit{Two-loop corrections}: The two-loop contribution is of order
        $g^2/(16\pi^2)^2 \cdot n \sim 0.01 \times 137 \approx 1.4$ for
        $g \sim 0.3$.  This accounts for roughly half the gap but has not been
        computed explicitly.
  \item \textit{Gauge-boson loops}: One-loop contributions from $\mathrm{SU}(2)$
        gauge bosons in the KK tower add $N_c b_0/(16\pi^2)$ corrections.
        Numerical estimate: $\delta B_\mathrm{gauge} \lesssim 1$.
  \item \textit{Non-perturbative threshold corrections}: At the
        compactification scale, non-perturbative contributions from Euclidean
        D-brane instantons (or analogous UBT objects) could provide a
        multiplicative factor $\approx 1.054$.
  \item \textit{Modular correction}: The modular index route (Section~\ref{sec:modular})
        gives the correct $B$ value but requires a non-perturbative input.
\end{enumerate}
\end{gapbox}

\subsection{Conclusion of Route 2}

\begin{resultbox}
Route 2 gives: $B_{\mathrm{KK+wind}} \approx 43.6$.
Status: \heur\ (closest perturbative estimate; short by factor 1.054).
Self-dual radius is assumed, not derived.  The gap $43.6 \to 46$ remains open.
\end{resultbox}

%────────────────────────────────────────────────────────────────────────────
\section{Route 3: Modular Index $\mu(\Gamma_0(p))/3$}
\label{sec:modular}
%────────────────────────────────────────────────────────────────────────────

\subsection{The modular invariant}

For a prime $p$, the index of the congruence subgroup $\Gamma_0(p)$ in
$\mathrm{SL}(2,\mathbb{Z})$ is (Diamond \& Shurman, Theorem 3.1.1):
\begin{equation}
  \mu(\Gamma_0(p)) = p + 1.
  \label{eq:modindex}
\end{equation}
This equals $|\mathbb{P}^1(\mathbb{F}_p)|$, the projective line over the finite
field $\mathbb{F}_p$.  The hyperbolic volume of the modular curve $X_0(p)$ is
\begin{equation}
  \mathrm{vol}(X_0(p)) = \frac{\pi}{3}\,\mu(\Gamma_0(p)) = \frac{\pi(p+1)}{3}.
  \label{eq:volX0}
\end{equation}
Dividing by $\pi$, the normalised hyperbolic area is
\begin{equation}
  \frac{\mathrm{vol}(X_0(p))}{\pi} = \frac{p+1}{3} = B(p),
  \label{eq:BmodIndex}
\end{equation}
which is precisely the phenomenological formula~\eqref{eq:Bformula}.

At $p = 137$:
\begin{equation}
  \frac{\mu(\Gamma_0(137))}{3} = \frac{138}{3} = 46.
  \label{eq:B137modular}
\end{equation}
The comparison with $B_{\mathrm{phenom}} = 46.298$ gives an error of $0.64\%$.
With the elliptic correction $\nu_2/4 = 0.5$ (from the two order-2 elliptic
points of $\Gamma_0(137)$, since $137 \equiv 1 \pmod{4}$):
\begin{equation}
  B_{\mathrm{mod+ell}} = \frac{\mu(\Gamma_0(p))}{3} + \frac{\nu_2}{4}
  = 46 + 0.5 = 46.5,
  \label{eq:BmodEll}
\end{equation}
reducing the error to $0.44\%$.

\subsection{Why $\mu(\Gamma_0(p))/3$ is structurally non-trivial}

The expression~\eqref{eq:BmodIndex} is:
\begin{itemize}
  \item A purely arithmetic invariant of $p$, computable without any physical input.
  \item Related to the normalised hyperbolic area of $X_0(p)$ by the formula
        $\mathrm{vol}(X_0(p))/(\pi/3) = \mu(\Gamma_0(p)) = p+1$.
  \item Equal to the degree of the modular polynomial $\Phi_p(X, Y)$ minus 1.
  \item Consistent with the projective geometry of $\mathbb{P}^1(\mathbb{F}_p)$:
        the cosets of $\Gamma_0(p) \backslash \mathrm{SL}(2,\mathbb{Z})$ biject with
        $\mathbb{P}^1(\mathbb{F}_p)$.
\end{itemize}

\subsection{Can $\mu(\Gamma_0(p))/3$ enter the effective action?}

\begin{gapbox}
\textbf{Gap 3 (modular index in $S[\Theta]$)}: The key open question is whether
the modular index $\mu(\Gamma_0(p)) = p+1$ can arise as a coefficient in the
one-loop effective action $\Gamma[\bar\Theta]$ evaluated on the winding-mode
background with level $n = p$.

\textbf{Mechanism candidates}:
\begin{enumerate}
  \item \textit{Hecke operator trace}: The one-loop effective action on the UBT
        torus $\Theta(q,\tau)$ can be written as a sum over Hecke orbits,
        $\mathrm{Tr}(T_p) = p + 1$ for the Hecke operator $T_p$ acting on the
        space of functions on $\mathbb{P}^1(\mathbb{F}_p)$.  If the UBT path
        integral measure weights the Hecke orbits uniformly, a factor $p+1$
        emerges naturally.
  \item \textit{Winding sector degeneracy}: In the winding sector at level $n$,
        the number of distinct coset representatives of $\Gamma_0(n)$ in
        $\mathrm{SL}(2,\mathbb{Z})$ is $\mu(\Gamma_0(n)) = n+1$ (for prime $n$).
        If this enumerates the winding vacua contributing to the loop sum,
        the effective coupling receives a factor $n+1$.
  \item \textit{Spectral geometry at $\tau = i$}: The Epstein zeta function for
        the square lattice $\mathbb{Z}[i]$ at the CM point $\tau = i$ couples
        the Hecke structure to the spectral determinant of the Laplacian (see
        \texttt{research\_tracks/T3\_ALPHA/chowla\_selberg\_b\_derivation.tex}).
        The Chowla--Selberg formula gives $B_{\mathrm{CS}} \approx 46.28$ close
        to the modular index value, but identifying the precise origin of $p+1$
        in this context remains an open calculation.
\end{enumerate}

\textbf{What is missing}: A first-principles computation of the UBT path
integral $\int \mathcal{D}\Theta\, e^{iS[\Theta]}$ in the winding sector at
level $n = p$ that returns a loop coefficient proportional to $\mu(\Gamma_0(p))$.
This requires either:
\begin{itemize}
  \item An exact calculation using the modular symmetry of $S[\Theta]$ (the
        action is invariant under $\mathrm{SL}(2,\mathbb{Z})$ transformations
        of $\tau$), or
  \item A non-perturbative saddle-point expansion that enumerates the
        $\mu(\Gamma_0(p))$ coset representatives explicitly.
\end{itemize}
\end{gapbox}

\subsection{Gauge invariance of the modular route}

\begin{proposition}[Gauge invariance of $\mu(\Gamma_0(p))/3$]
  The modular index $\mu(\Gamma_0(p)) = p + 1$ is invariant under the SM gauge
  group $\mathrm{SU}(3)\times\mathrm{SU}(2)\times\mathrm{U}(1)$: it is an
  arithmetic invariant of the prime $p$ and does not depend on any gauge field
  configuration.
\end{proposition}

\begin{proof}
  The index $\mu(\Gamma_0(p))$ counts cosets in $\mathrm{SL}(2,\mathbb{Z})$ and
  is determined solely by $p$.  The modular group $\mathrm{SL}(2,\mathbb{Z})$
  acts on the complex-time torus $\tau = t + i\psi$; it commutes with the SM
  gauge transformations acting on the biquaternion field components.  Hence
  $\mu(\Gamma_0(p))$ is SM-gauge-invariant. \qed
\end{proof}

\subsection{Regulator independence of the modular route}

The modular index is a topological (arithmetic) invariant, independent of any
UV regulator.  The Hecke degeneracy $p+1$ of the winding sector is an exact
result of the discrete geometry of $\mathbb{P}^1(\mathbb{F}_p)$ and is not
affected by renormalisation scheme choices.

\subsection{Conclusion of Route 3}

\begin{resultbox}
Route 3 gives: $B_{\mathrm{mod}} = (p+1)/3 = 46$ at $p = 137$ (error 0.64\%).
Status: \obs\ (strongly motivated coincidence; not yet derived from $S[\Theta]$).
The modular index is gauge-invariant and regulator-independent.  The missing
step is a derivation showing that $\mu(\Gamma_0(n^*))$ enters the one-loop
effective action coefficient.
\end{resultbox}

%────────────────────────────────────────────────────────────────────────────
\section{Route 4: Fixed-Point Bootstrap Self-Consistency}
\label{sec:bootstrap}
%────────────────────────────────────────────────────────────────────────────

\subsection{The bootstrap system}

Combining the prime-attractor condition~\eqref{eq:fixedpoint} with the
modular-index ansatz~\eqref{eq:Bformula} gives the system:
\begin{align}
  2n &= B\bigl(\ln n + 1\bigr), \label{eq:bs1} \\
  B  &= \frac{n+1}{3}.           \label{eq:bs2}
\end{align}
Substituting~\eqref{eq:bs2} into~\eqref{eq:bs1}:
\begin{equation}
  6n = (n+1)(\ln n + 1).
  \label{eq:bootstrap}
\end{equation}

\subsection{Numerical verification}

Define $f(n) = (n+1)(\ln n + 1) - 6n$.  Then:

\begin{center}
\begin{tabular}{cccc}
\toprule
$n$ & $f(n)$ & Sign & Nearest prime \\
\midrule
100 & $(101)(5.605) - 600 = 566.1 - 600 = -33.9$ & $-$ & --- \\
130 & $(131)(5.567) - 780 = 729.3 - 780 = -50.7$ & $-$ & --- \\
137 & $(138)(5.919) - 822 = 816.8 - 822 = -5.2$ & $-$ & $\mathbf{137}$ \\
138 & $(139)(5.927) - 828 = 823.8 - 828 = -4.2$ & $-$ & 139 \\
139 & $(140)(5.934) - 834 = 830.8 - 834 = -3.2$ & $-$ & $\mathbf{139}$ \\
150 & $(151)(6.008) - 900 = 907.2 - 900 = +7.2$ & $+$ & 151 \\
\bottomrule
\end{tabular}
\end{center}

The root of $f(n) = 0$ lies in $[139, 150]$; the nearest prime is
$n = 139$ on one side and $n = 137$ on the other.  The exact solution
$n_{\mathrm{root}} \approx 141$ is not prime.

\begin{insightbox}
\textbf{Interpretation}: The bootstrap system~\eqref{eq:bs1}--\eqref{eq:bs2}
does not have an exact prime solution.  The prime $n = 137$ is the largest
prime for which $f(p) < 0$, i.e.\ the modular formula slightly undershoots the
prime-attractor condition (by 0.64\%).  The prime $n = 139$ also satisfies the
system approximately ($f(139) \approx -3.2$, error 0.39\%).  Neither prime is
an exact solution; both are approximate fixed points.

This is consistent with the modular interpretation: the exact solution requires
the elliptic correction $\varepsilon = \nu_2/4 \approx 0.5$ to match
$B_{\mathrm{phenom}} \approx 46.298$ to better accuracy.
\end{insightbox}

\subsection{Bootstrap as a self-consistency check, not a derivation}

The bootstrap~\eqref{eq:bootstrap} confirms that the system~\eqref{eq:Bformula} is
\emph{approximately self-consistent} near $n = 137$ and $n = 139$.  However,
it \emph{does not} derive~\eqref{eq:Bformula} from $S[\Theta]$: it assumes
$B = (n+1)/3$ and verifies consistency.  A derivation must produce
$B = (n+1)/3$ as an output of $S[\Theta]$, not as an input.

\begin{resultbox}
Route 4 gives: Self-consistency confirmed near $n = 137$.
Status: \cond\ (the ansatz is self-consistent; it is not a derivation).
No new information about the origin of $B(p) = (p+1)/3$ is provided.
\end{resultbox}

%────────────────────────────────────────────────────────────────────────────
\section{Summary of All Routes}
\label{sec:summary}
%────────────────────────────────────────────────────────────────────────────

\begin{longtable}{p{3.5cm}p{2.5cm}p{3.5cm}p{4cm}}
\toprule
\textbf{Route} & \textbf{$B$ estimate} & \textbf{Status} & \textbf{Key assumption/gap} \\
\midrule
One-loop KK ($d=1$) & 21.8 & \heur & Self-dual radius; $n$-dependent \\
KK + winding & 43.6 & \heur & T-duality at $R=1$ (not derived) \\
Two-loop correction & $+1$--$2$ & \heur & Diagrams not computed \\
\textbf{KK+wind+2-loop} & $\approx 44$--$46$ & \heur & Factor 1.054 unexplained \\
Modular index / $\Gamma_0(p)$ & 46.0 & \obs & $\mu(\Gamma_0(p))$ not in $S[\Theta]$ \\
Mod.\ + elliptic correction & 46.5 & \obs & Same; $\varepsilon$ origin open \\
Fixed-point bootstrap & self-consistent & \cond & Ansatz, not derivation \\
\midrule
\textbf{Target} & \textbf{46.298} & --- & --- \\
\bottomrule
\caption{Summary of $B$ estimates from all routes.}
\label{tab:summary}
\end{longtable}

%────────────────────────────────────────────────────────────────────────────
\section{Gap Matrix}
\label{sec:gapmatrix}
%────────────────────────────────────────────────────────────────────────────

\begin{center}
\begin{tabular}{p{4.5cm}lp{5.5cm}}
\toprule
\textbf{Claim} & \textbf{Status} & \textbf{Gap / remark} \\
\midrule
$n^2$ term from KK & \proved & --- \\
$n^2 \ln n$ shape from one-loop & \proved & --- \\
Regulator independence & \proved & Standard $d=1$ dim.\ reg. \\
Gauge invariance of $V_{\mathrm{eff}}$ & \proved & $n$ is $\mathrm{U}(1)_\psi$ charge \\
$B$ constant (not $n$-dependent) & \heur & Requires fixing $n = n^*$ \\
$B \approx 21.8$ (one-loop) & \heur & Self-dual radius assumed \\
$B \approx 43.6$ (KK+winding) & \heur & T-duality factor 2 \\
$B = 46$ (modular) & \obs & $\mu(\Gamma_0(p))$ not in $S[\Theta]$ \\
Elliptic correction $\varepsilon = 0.5$ & \obs & $\nu_2/4$ numerics only \\
$B(p) = (p+1)/3$ derived & \openp & \textbf{Gap G137-B: OPEN} \\
Bootstrap self-consistency & \cond & Verified, not derived \\
\bottomrule
\end{tabular}
\end{center}

%────────────────────────────────────────────────────────────────────────────
\section{Verdict}
\label{sec:verdict}
%────────────────────────────────────────────────────────────────────────────

\begin{tcolorbox}[colback=red!5!white, colframe=red!65!black, title=\textbf{Verdict on Gap G137-B}]
\textbf{$B(p) = (p+1)/3$ CANNOT currently be derived from the UBT action
$S[\Theta]$.}

\medskip

The formula must remain a \textbf{CONDITIONAL modular ansatz} until one of the
following is achieved:
\begin{enumerate}
  \item A computation of the UBT path integral in the winding sector at level
        $n = p$ that returns the Hecke-degeneracy factor $\mu(\Gamma_0(p)) = p+1$
        as the effective coupling.
  \item An explicit two-loop + non-perturbative calculation that closes the
        $43.6 \to 46$ gap without modular input.
  \item A proof that the modular symmetry of $S[\Theta]$ forces the one-loop
        coefficient to equal $\mathrm{vol}(X_0(n^*))/\pi$.
\end{enumerate}

\medskip
\noindent\textbf{Current best estimate}: $B \approx 43.6$--$44$ from the
RG/KK+winding route.  The modular route gives $B = 46$ exactly at $p = 137$
but is not derived from $S[\Theta]$.

\medskip
\noindent\textbf{Experimental input}: None.  The value $\alpha = 1/137$ was
not used in any step.
\end{tcolorbox}

\subsection{Priority recommendation}

Given the difficulty of the remaining gap (27+ failed approaches, see
\texttt{reports/alpha\_missing\_lemma.md}) and the time-box constraint, the
recommended path is:

\begin{enumerate}
  \item \textbf{Time-box 4 weeks}: Attempt the Hecke-trace route (Section~\ref{sec:modular},
        Gap 3, candidate mechanism 1).  This is the most structurally motivated
        attack.
  \item \textbf{If unsuccessful}: Publish the integer-137 result as
        \textbf{CONDITIONAL} on Gap G137-B, with $B(p) = (p+1)/3$ stated
        explicitly as a modular ansatz supported by three coincidences
        (modular index, fixed-point consistency, and the Chowla--Selberg
        numerical match).
  \item \textbf{Refocus resources} on T1\_GR (general relativity sector) and
        T2\_GAUGE (gauge sector), where closure is closer.
\end{enumerate}

%────────────────────────────────────────────────────────────────────────────
\section{Falsification Conditions}
\label{sec:falsification}
%────────────────────────────────────────────────────────────────────────────

\begin{center}
\begin{tabular}{p{6cm}p{7.5cm}}
\toprule
\textbf{Condition} & \textbf{What it falsifies} \\
\midrule
Two-loop calculation gives $B_\mathrm{2-loop} < 44$ & RG+perturbative origin of $B = 46$ \\
$R_\psi \neq 1$ from UBT moduli derivation & Self-dual radius assumption \\
$B(p)$ not monotone in $p$ & Prime-stability formula $B = (p+1)/3$ \\
Modular symmetry of $S[\Theta]$ is broken & Hecke-trace mechanism \\
$n^*(B_\mathrm{phenom}) \neq 137$ at higher order & Entire prime-attractor programme \\
\bottomrule
\end{tabular}
\end{center}

%────────────────────────────────────────────────────────────────────────────
\begin{thebibliography}{9}
\bibitem{DSh05} F.\ Diamond and J.\ Shurman, \emph{A First Course in Modular
  Forms}, Springer, 2005.
\bibitem{ZinnJustin} J.\ Zinn-Justin, \emph{Quantum Field Theory and Critical
  Phenomena}, 4th ed., Oxford, 2002.
\bibitem{CW73} S.\ Coleman and E.\ Weinberg, ``Radiative corrections as the
  origin of spontaneous symmetry breaking,'' \emph{Phys.\ Rev.\ D}
  \textbf{7} (1973) 1888.
\bibitem{CS49} S.\ Chowla and A.\ Selberg, ``On Epstein's zeta function,''
  \emph{J.\ reine angew.\ Math.} \textbf{227} (1967) 86.
\bibitem{Tong} D.\ Tong, \emph{Lectures on String Theory}, Cambridge, 2009.
\bibitem{rg_nlogn} Ing.\ D.\ Jaro\v{s},
  \texttt{research\_tracks/rg\_nlogn/full\_rg\_derivation.tex}, 2026.
\bibitem{gamma0} Ing.\ D.\ Jaro\v{s},
  \texttt{reports/gamma0\_137\_invariants.md}, 2026.
\bibitem{cs_route} Ing.\ D.\ Jaro\v{s},
  \texttt{research\_tracks/T3\_ALPHA/chowla\_selberg\_b\_derivation.tex}, 2026.
\bibitem{missing} Ing.\ D.\ Jaro\v{s},
  \texttt{reports/alpha\_missing\_lemma.md}, 2026.
\end{thebibliography}

\end{document}
